\documentclass[12pt]{article}
\usepackage[utf8]{inputenc}
\usepackage[T1]{fontenc}
\usepackage[brazil]{babel}
\usepackage{graphicx} % Required for inserting images
\usepackage[a4paper,left=3cm,right=2cm,top=3cm,bottom=2cm]{geometry} % Margens ABNT
\usepackage{setspace} % Espaçamento entre linhas
\onehalfspacing % Espaçamento 1,5 (padrão ABNT)
\usepackage[alf,num=false,abnt-etal-cite=2,abnt-etal-list=0,abnt-etal-text=it]{abntex2cite}

% Dados do trabalho
\newcommand{\nomealuno}{Bruno Martins Valério Bomfim}
\newcommand{\titulotrabalho}{Uma proposta de implementação de ambiente virtual para prover recursos virtuais para disciplinas de computação}
\newcommand{\nomeorientador}{Prof. Dr. Fernando William Cruz}
\newcommand{\instituicao}{Faculdade de Ciências e Tecnologias em Engenharia/Universidade de Brasília — FCTE/UnB}
\newcommand{\programa}{Programa de Iniciação Científica}
\newcommand{\anorealizacao}{2026}

\begin{document}

% ---------- TÍTULO ----------
\begin{center}
    {\large \nomealuno}

    \vspace{0.3cm}

    {\Large \bfseries \titulotrabalho}

    \vspace{0.3cm}

    {\large \nomeorientador}

    \vspace{0.3cm}

    {\large \instituicao}

    \vspace{0.3cm}

    {\large \programa}

    \vspace{0.3cm}

    {\large \anorealizacao}
\end{center}

\vspace{1cm}

% ---------- RESUMO ----------
\section*{Resumo}
Texto do resumo (até 250 palavras).

\vspace{0.5cm}
\noindent \textbf{Palavras-chave:} palavra1, palavra2, palavra3.

\newpage

% ---------- INTRODUÇÃO ----------
\section{Introdução}
\label{intro}

Contextualize o tema, apresente os objetivos e a justificativa do projeto.


% ---------- REFERENCIAL TEÓRICO ----------
\section{Referencial Teórico}

% ---------- METODOLOGIA ----------
\section{Metodologia}

Esta seção descreve os métodos que direcionaram as atividades realizadas neste projeto de iniciação científica. Apresenta-se também a classificação desta pesquisa e o respectivo levantamento bibliográfico, além de detalhamentos sobre a prova de conceito e o cronograma de atividades.

Descreva os métodos e técnicas utilizados na realização do trabalho.

\subsection{Classificação da Pesquisa}
\label{test}

\citeonline[p.~33]{gerhardt_silveira} definem a pesquisa científica como \lq\lq o resultado de um inquérito ou exame minucioso, realizado com o objetivo de resolver um problema, recorrendo a procedimentos científicos\rq\rq. Seguindo a taxonomia proposta por estas autoras, é apresentada abaixo a classificação desta pesquisa quanto à \textbf{abordagem}, \textbf{natureza}, \textbf{objetivo} e \textbf{procedimentos}.

\subsubsection{Abordagem}

Esta pesquisa possui abordagem \textbf{qualitativa e quantitativa} \cite[p. ~33-36]{gerhardt_silveira}. O caráter \textbf{qualitativo} está relacionado à análise criteriosa de técnicas existentes na literatura científica para a resolução do problema proposto na seção \ref{intro}. De forma complementar, o aspecto \textbf{quantitativo} está relacionado à coleta e estudo de métricas concernentes à performance da solução proposta (veja \ref{levantamento}).

\subsubsection{Natureza}

Esta pesquisa possui natureza \textbf{aplicada} \cite[p. ~36-37]{gerhardt_silveira}, haja vista que almeja engendrar conhecimentos para a solução prática de um problema específico, a saber, o problema de compartilhamento dos laboratórios da FCTE/UnB.

\subsubsection{Objetivos}

Os objetivos deste estudo são \textbf{exploratórios} \cite[p.~37-38]{gerhardt_silveira}, envolvendo, assim, levantamento bibliográfico, e o estudo de caso do laboratório LDS da FCTE.

\subsubsection{Procedimentos}

Adotou-se a \textbf{pesquisa bibliográfica} \cite[p.~39]{gerhardt_silveira} cujo objetivo concentrou-se em identificar soluções propostas pela literatura científica para a implementação de gêmeos digitais em ambiente de realidade virtual ou aumentada. Ademais, a \textbf{pesquisa experimental} \cite[p.~38]{gerhardt_silveira} buscou construir uma solução para virtualização do laboratório LDS e a construção de uma prova de conceito no ambiente Decentraland.

\subsection{Levantamento Bibliográfico}
\label{levantamento}

\subsubsection{Strings de Busca}

\subsubsection{Critérios de Seleção}

\subsection{Prova de Conceito}

\subsection{Cronograma}

% ---------- RESULTADOS E DISCUSSÃO ----------
\section{Resultados e Discussão}

Apresente os resultados obtidos e discuta-os criticamente.


% ---------- CONSIDERAÇÕES FINAIS / CONCLUSÃO ----------
\section{Considerações Finais / Conclusão}

Síntese dos achados, contribuições, limitações e sugestões futuras.

% ---------- REFERÊNCIAS ----------
\newpage
\bibliography{references} 

\end{document}